
\section{Introduction}

In this paper we present a unified framework for the valuation of caps, floors and swaptions. These instruments are the most common derivative securities which are traded in a fixed income desk of a financial institution (see e.g. \cite{brigo06}). Practitioners usually price these products by relying on a Black-Scholes like formula, which was first presented in \cite{article_Black1976}. The market convention of pricing caps and swaptions using the Black formula is based on an application of the \cite{bla73} formula for stock options by assuming that the underlying interest rates are lognormally distributed. Remarkably, the use of this kind of formulae had no theoretical justification, since they involved a procedure in which the discount factor and the Libor rates were assumed to be independent so as to write the pricing formula as a product of a bond price and the expected payoff. The systematic use of this market practice ignited the interest of academics aiming at providing a coherent theoretical background.\\

In a series of articles, \cite{article_MiltersenSandmannSondermann1997}, \cite{article_BraceGatarekMusiela1997}, \cite{article_Jamshidian1997} and \cite{article_MusielaRutkowsky1997}, provided these theoretical foundations, introducing the Libor and Swap Market Model. Following these papers a stream of contributions appeared, trying to extend the basic model to the case where the volatility of the underlying factor is stochastic. The most famous proposals on this side can be found e.g. in \cite{article_AndersenBrotherton-Ratcliffe2001}, \cite{article_WuZhang2006},  \cite{article_JoshiRebonato2003}, \cite{andand2002}, \cite{pit05}, \cite{pit05b}. Other approaches explored different dynamics for the driving process with respect to the CEV or displaced-diffusion considered before for the Libor rate: for example \cite{glass03}, \cite{article_EberleinOzkan2005} introduced jump and more general L�vy processes, allowing for discontinuous sample paths of the driving process. Another interesting approach is the one of \cite{brigo03} based on a mixture of lognormals.\\
% A common pitfall consists in the problematic form of the variance-covariance matrix of the yields, which makes it difficult to perform a PCA analysis.\\

A typical problem in the previous approaches is that once the closed form solution for cap prices is found, to obtain an analogous result for swaptions it is customary to assume that the underlying (which is a coupon bond) behaves like a scalar process (typically again geometric Brownian motion). This results in inconsistencies between the so-called \textit{Libor} and \textit{Swap Market Models}. Even more important, by assuming that the coupon bond is driven by a scalar process, we do not take into account the correlation effects among the different coupons, a key feature of a swaption which may be viewed also as a correlation product. This last remark is of paramount importance for practitioners (see e.g. the introduction of \cite{article_CollinDufresneGoldstein2002}).\\



In this paper we consider a new approach based on the stochastic discount factor methodology, where instead of modeling directly the Libor rate, one concentrates on quotients of traded assets (i.e. bonds).  It has been first introduced by \cite{article_Constantinides92} and then developed by \cite{gou11} in a spot interest rate framework and by \cite{article_KRPT2009} in a Libor perspective. In this latter work they use affine processes on the state space $\mathbb{R}_{\geq0}^{d}$ as driving processes and provide a full characterization of the model, which allows them to provide closed form solutions for caps and swaptions up to Fourier integrals. This approach is very interesting and easily overcomes many difficulties which are to be faced in the computation of Radon-Nikodym derivatives.\\

We provide an extension of this approach, by considering affine processes on the state space $S_{d}^{++}$, the set of positive definite symmetric matrices. This state space may seem awkward at first sight, but the processes belonging to this family admit a characterization in terms of ODE's which resembles the one found for standard affine models, an example being given by the famous \cite{article_DuffieKan} model. In fact, in \cite{article_Cuchiero} the authors extend to the set $S_{d}^{+}$ (the set of positive semidefinite symmetric matrices) the classification of affine processes performed by \cite{article_DFS} for the state space $\mathbb{R}_{\geq0}^{d}\times \mathbb{R}^{n-d}$ introduced by \cite{article_DuffieKan}. What is more, the state space $S_{d}^{++}$ leads to stochastic factors which are non trivially correlated.  The most famous example of process defined in the set $S_{d}^{++}$ is the Wishart process, originally defined by \cite{article_Bru}, introduced in finance by \cite{article_gou02} and then extensively applied in \cite{article_gou03}, \cite{gou11}, \cite{article_DaFonseca1}, \cite{article_DaFonseca2}, \cite{article_DaFonseca3}, \cite{dafgra11}, among others.\\ 

The interesting feature of our framework is the possibility to obtain semi-closed form solutions for the pricing of swaptions in a multifactor setting, which is a well known challenging problem. In fact the exercise probability involves a multi-dimensional inequality. There have been many approaches to simplify the problem: for example, \cite{sinum2002} suggest an approximation of the exercise boundary
with a linear function of the state variables. However, the most efficient approach seems to be the one of \cite{article_CollinDufresneGoldstein2002} which heavily uses the affine structure of the model and is based on the Edgeworth expansion for the characteristic function in terms of the cumulants.  Since the cumulants decay very quickly the Edgeworth expansion for the exercise probability turns out to be very accurate and fast.\\


The paper is organized as follows. In Section 2 we introduce our framework by recalling some useful definitions and results on affine processes. Section 3 investigates the case of the state space $S_{d}^{++}$ and presents the technical results. In Section 4 we focus on the pricing problem of the relevant derivatives. Caps and Floors are briefly treated since their pricing is now quite standard within the FFT methodology, while we devote more attention to the pricing of swaptions by adopting the approach of \cite{article_CollinDufresneGoldstein2002}. Section 5 illustrates the flexibility of our framework through a numerical exercise. Section 6 concludes the paper, and we gather in the technical Appendices proofs and some remarks useful for implementation.\\ 






